\lecture{8.1}{Longitudinal Causal Questions}
\label{lec:8-1}

\subsection{Multi-Timepoint Treatments}

Many important interventions are not single decisions. A patient may take a
medication every day, receive chemotherapy at several cycles, return for repeated
injections, or decide whether to continue a behavioral program at each visit. In
such settings the causal question concerns a \emph{sequence} of treatment
decisions, not only one baseline exposure.

Let time be indexed by $k=0,1,\ldots,K$. At time $k$, an individual has a
recorded history, receives or does not receive treatment, and then continues to
the next time point. The final outcome $Y$ is observed after the last treatment
decision.

\begin{figure}[H]
\centering
\begin{tikzpicture}[
  node distance=1.05cm,
  box/.style={draw, rounded corners=2pt, minimum width=1.05cm, minimum height=.72cm},
  arr/.style={-{Latex[length=2.2mm]}, thick}
]
\node[box] (L0) {$L_0$};
\node[box, right=of L0] (A0) {$A_0$};
\node[box, right=of A0] (L1) {$L_1$};
\node[box, right=of L1] (A1) {$A_1$};
\node[right=of A1] (dots) {$\cdots$};
\node[box, right=of dots] (LK) {$L_K$};
\node[box, right=of LK] (AK) {$A_K$};
\node[box, right=of AK] (Y) {$Y$};
\draw[arr] (L0) -- (A0);
\draw[arr] (A0) -- (L1);
\draw[arr] (L1) -- (A1);
\draw[arr] (A1) -- (dots);
\draw[arr] (dots) -- (LK);
\draw[arr] (LK) -- (AK);
\draw[arr] (AK) -- (Y);
\draw[arr, bend left=25] (A0) to (Y);
\draw[arr, bend left=22] (A1) to (Y);
\draw[arr, bend left=18] (LK) to (Y);
\end{tikzpicture}
\caption{A generic longitudinal treatment process. Covariates and treatments alternate over time before the final outcome.}
\label{fig:ch8-longitudinal-timeline}
\end{figure}

\begin{keyidea}
A longitudinal causal question asks what would happen if the entire treatment
history were set by an intervention. The unit of intervention is the treatment
regime $\abar=(a_0,\ldots,a_K)$, not a single treatment indicator.
\end{keyidea}

\subsection{Longitudinal Data Structure}

For one individual, write the observed data as
\[
  O=(L_0,A_0,L_1,A_1,\ldots,L_K,A_K,Y).
\]
For $n$ independent individuals, the sample is
\[
  O_1,\ldots,O_n \iid P_0,
\]
where $P_0$ denotes the true observed-data distribution. The variables have the
following roles:
\[
  L_k = \text{covariates recorded at time } k,\qquad
  A_k = \text{treatment at time } k,\qquad
  Y = \text{final outcome}.
\]
The covariate $L_k$ may include new measurements, updated disease severity,
adverse events, adherence information, symptoms, laboratory values, and summaries
of earlier history. It is often vector-valued.

\begin{notation}
We use overbars for histories:
\[
  \Lbar_k=(L_0,\ldots,L_k),\qquad
  \Abar_k=(A_0,\ldots,A_k).
\]
For fixed treatment values,
\[
  \lbar_k=(l_0,\ldots,l_k),\qquad
  \abar_k=(a_0,\ldots,a_k).
\]
When the full treatment history is meant, $\Abar=\Abar_K$ and
$\abar=\abar_K$.
\end{notation}

The timing convention matters. In these notes, $L_k$ is measured before $A_k$.
Thus $A_k$ may depend on $\Lbar_k$ and $\Abar_{k-1}$, but not on future
covariates or future outcomes. This convention is not merely bookkeeping: it
determines what information can be used by a treatment rule at time $k$.

\subsection{$L_k$, $A_k$, and $Y$}

The variable $L_k$ can be both a confounder and an intermediate. For example,
baseline treatment $A_0$ may affect a later biomarker $L_1$, and that biomarker
may affect both later treatment $A_1$ and the outcome $Y$. This feedback is a
central reason why longitudinal causal inference requires special methods.

For two treatment times, the observed data are
\[
  O=(L_0,A_0,L_1,A_1,Y).
\]
A typical causal ordering is
\[
  L_0 \rightarrow A_0 \rightarrow L_1 \rightarrow A_1 \rightarrow Y,
\]
with additional arrows from earlier variables to later variables. The later
covariate $L_1$ is post-treatment with respect to $A_0$, but pre-treatment with
respect to $A_1$.

\subsection{Treatment Regimes}

\begin{definition}[Static treatment regime]
A static treatment regime is a fixed vector
\[
  \abar=(a_0,\ldots,a_K)
\]
that assigns the same treatment value at each time point to every individual,
regardless of the individual's evolving covariate history.
\end{definition}

Examples include ``always treat,'' represented by $(1,1,\ldots,1)$, and
``never treat,'' represented by $(0,0,\ldots,0)$ for a binary treatment.
Other static regimes are possible, such as treating only during the first two
visits and then stopping.

\begin{definition}[Dynamic treatment regime]
A dynamic treatment regime is a sequence of decision rules
\[
  d=(d_0,\ldots,d_K),
\]
where $d_k$ maps the information available at time $k$ to a treatment value:
\[
  A_k=d_k(\Lbar_k,\Abar_{k-1}).
\]
\end{definition}

Dynamic regimes are discussed in detail in \cref{lec:8-11}. For now, the key
point is that longitudinal notation can handle both fixed treatment profiles and
rules that respond to patient history.

\subsection{Counterfactual Outcomes $Y(a_0,\ldots,a_K)$}

\begin{definition}[Longitudinal counterfactual outcome]
For a fixed treatment regime $\abar=(a_0,\ldots,a_K)$, the counterfactual
outcome $Y(\abar)$ is the outcome that would be observed for an individual if,
possibly contrary to fact, treatment were set to $a_k$ at every time
$k=0,\ldots,K$.
\end{definition}

For example, $Y(1,1)$ is the outcome that would be observed under treatment at
both time 0 and time 1. If the observed data contain $A_0=1$ and $A_1=0$, then
$Y(1,1)$ is not directly observed for that individual, even though $Y(1,0)$ may
be linked to the observed outcome by consistency.

Common population-level causal parameters include
\[
  \psi(\abar)=\E\{Y(\abar)\},
\]
and contrasts such as
\[
  \E\{Y(1,\ldots,1)\}-\E\{Y(0,\ldots,0)\}.
\]
For a binary outcome, this contrast is a risk difference comparing two entire
treatment strategies.

\subsection{Clinical Examples}

\begin{example}[Weekly osteoporosis therapy]
Suppose $A_k$ indicates whether a patient took alendronate during week $k$,
where $k=0,\ldots,52$. Let $L_k$ include age, prior fracture status, recent bone
density measurements, thyroid hormone use, side effects, and other clinical
information available at week $k$. Let $Y$ indicate whether hip fracture occurs
within one year. A static contrast of interest is
\[
  \E\{Y(1,1,\ldots,1)\}-\E\{Y(0,0,\ldots,0)\},
\]
the difference in one-year fracture risk under weekly therapy for the whole year
versus no therapy during the year.
\end{example}

\begin{example}[Cancer treatment continuation]
At each chemotherapy cycle, $A_k$ may indicate whether treatment is continued.
The covariate $L_k$ may include tumor response, toxicity, blood counts, and
functional status. Treatment discontinuation may be caused by lack of response,
and lack of response also predicts mortality. Thus later treatment decisions are
confounded by earlier treatment-induced covariates.
\end{example}

\begin{example}[Smoking cessation]
Let $A_k$ indicate whether an individual smokes during month $k$. Respiratory
symptoms recorded in $L_k$ may lead a person to quit smoking and may also be
early signs of disease. A comparison of sustained smoking versus sustained
cessation must handle the fact that symptoms are both affected by prior smoking
and predictive of later smoking and future outcomes.
\end{example}
